\documentclass[10pt,a4paper]{amsart}
\usepackage[margin=19mm]{geometry}
\usepackage{amsmath,amssymb,mathtools}
\usepackage[scr=rsfs,cal=euler]{mathalfa}
\usepackage{xfrac}
\usepackage[unicode,hidelinks,bookmarks=false]{hyperref}
\newcommand{\ie}{i.e.\ }
\newcommand{\cf}{cf.\ }
\newcommand{\resp}{respectively\ }
\newcommand{\Subsection}{\subsection}
\providecommand{\nobreakdash}{\nobreak\hbox{-}}
\setlength{\emergencystretch}{2em}
\multlinegap0pt
\usepackage{bm}
\usepackage{bbm}
\usepackage{dsfont}
\usepackage{xspace}
\usepackage{cancel}
\usepackage{mathtools}
\usepackage{amsthm}
\makeatletter
\@ifundefined{thm}{\newtheorem{thm}{Thm}}{}
\@ifundefined{lem}{\newtheorem{lem}[thm]{Lemma}}{}
\makeatother
\def\cC{\mathcal C}
\title[Configurations in the Euclidean plane associated to a system]{Configurations in the Euclidean plane associated to a system of equations}
\author{Francesco Colangelo}
\date{Source: arXiv:2506.14773v1 (CC BY 4.0)}
\begin{document}
\maketitle

\noindent\emph{Source and licence.} Configurations in the Euclidean plane associated to a system of equations. Version arXiv:2506.14773v1 (CC BY 4.0), distributed under the Creative Commons Attribution 4.0 International licence, as recorded in the arXiv licence metadata for this exact version and shown on the abstract page https://arxiv.org/abs/2506.14773v1. Full rights notice: licenses/arxiv-2506.14773-CC-BY-4.0.txt.\par\smallskip

\noindent\textit{Editorial scope.} Counted unit: the Lem whose source label is \texttt{lemH071224} in the article (source0.tex lines 357--359, source label \texttt{lemH071224}), delivered together with the article's own complete original proof of it (source0.tex lines 360--421, one \texttt{proof} environment of 62 lines). No mathematics is written, repaired or re-derived: statement and proof are the article's own text line for line. Required context retained uncounted: lem20apr (lines 303--305); lem17apr (lines 316--318). Every definition, display and label of those lines is retained unchanged. Omitted from this package and not redistributed: 1--302, 306--315, 319--356, 422--624. Nothing inside a retained excerpt is omitted or altered: every retained body is delivered exactly as the article's lines are written, so no omission receipt of any kind accompanies this package. Citations: the retained excerpts carry no citation command at all, so no bibliography entry of the article is transcribed and no reference is redistributed. Reference adaptation: none was needed and none was made; every reference inside the delivered unit resolves inside it, so no unresolved reference or citation is produced. Printed slips kept literal: no printed word, symbol, hypothesis or number of the counted statement or of its proof was altered. Layout and class: the article's own class is \texttt{amsart} with packages including amssymb, amsmath, color; the wrapper is the lane's pinned \texttt{amsart} editorial wrapper at 10pt on A4, which restates the theorem-like environments the retained text needs and adds the article's own macro definitions that the retained text needs (\texttt{\textbackslash cC}), with extra wrapper packages (bm, bbm, dsfont, xspace, cancel, mathtools). The delivered numbering is the unit's own. Assets: no retained span contains a figure environment, a graphics-inclusion command, a PDF-page inclusion, a table or a TikZ picture; no image file is copied into this package, the package declares no asset dependency and no image file is redistributed with it. Licence: version arXiv:2506.14773v1 of this article is distributed under the Creative Commons Attribution 4.0 International licence, as recorded in the arXiv licence metadata for this exact version and shown on the abstract page https://arxiv.org/abs/2506.14773v1. A full rights notice is delivered at licenses/arxiv-2506.14773-CC-BY-4.0.txt. Characters: every retained excerpt is pure ASCII, so the delivered engine, XeLaTeX with its default Latin Modern text font, reports no missing character.
\begin{lem}
\label{lem20apr} For $T\in\{A,B,C,D\},$ neither $\|Y-T\|^2$ nor $k_T\|Y-T\|^2-1$ vanishes at each branch of $\cC.$
\end{lem}

\begin{lem}
\label{lem17apr} If the branch $\gamma$ of $\cC$ is centered at an affine point, then $\gamma$ is not a pole of $\Gamma_T(Y)$ with $T\in \{A,B,C,D\}$.
\end{lem}

\begin{lem}
\label{lemH071224} If the branch $\gamma$ is centered at a point of\, $\cC$ at infinity, then $\gamma$ is not a pole of $\Gamma_T(Y)$ for any $T\in \{A,B,C,D\}$.
\end{lem}

\begin{proof} Before Lemma \ref{lem20apr}, we have pointed out that $\cC$ has two points at infinity, namely $P^+$ and $P^-$. Let $P$ any of them. Then $P=(\eta_1:\eta_2:0)$ with $\eta_1^2+\eta_2^2=0$. Let $\gamma$ be a branch of $\cC$ centered at $P$ with a primitive branch representation
$(x(\tau):y(\tau):1)$ where
$$x(\tau)=\frac{\eta_1+\varphi_1(\tau)}{\alpha(\tau)},\quad y(\tau)=\frac{\eta_2+\varphi_2(\tau)}{\alpha(\tau)}$$
with $\alpha(\tau)=\tau^{m} \beta(\tau)$ with $m\ge 1$ and $\beta(\tau)\in \mathbb{C}[[\tau]],\  \beta(0)\neq 0$.

Suppose that  $\gamma$ is a pole of $\Gamma_A(Y)$.  Then $\gamma$ is also a pole of $\Gamma_A(Y)+\|A\|^2$, and hence $\gamma$ is a zero of $(\Gamma_A(Y)+\|A\|^2)^{-1}$, that is, $\gamma$ is a zero of
\begin{equation}
\label{eqF17apr}
k_A-\frac{1}{\|Y+A\|^2}.
\end{equation}
 By a straightforward computation,
$$
\|Y+A\|^2=(x(\tau)+a_1)^2+(y(\tau)+a_2)^2=
\frac{\sum\limits_{i=1}^2(\eta_i+\varphi_i(\tau)+a_i\alpha(\tau))^2}{\alpha(\tau)^2}=
$$



$$\frac{\sum\limits_{i=1}^2 \eta_i^2+2\sum\limits_{i=1}^2 \eta_ia_i\alpha(\tau)+2\sum\limits_{i=1}^2 a_i\alpha(\tau)\varphi_i(\tau)+\|A\|^2\alpha(\tau)^2+2\sum\limits_{i=1}^2 \eta_i\varphi_i(\tau)
+\sum\limits_{i=1}^2 \varphi_i^2(\tau)}{\alpha(\tau)^2}.
$$
By expanding (\ref{eqF17apr}),
$$k_A- \frac{ 1}{W_1(\tau)+W_2(\tau)+W_3(\tau)+Z_1(\tau)+Z_2(\tau)}$$
where
$$
W_1(\tau)=2\Big(\sum\limits_{i=1}^2 \eta_ia_i\Big)\tau^{-m}\beta(\tau)^{-1},\,
W_2(\tau)=2\Big(\sum\limits_{i=1}^2 a_i\varphi_i(\tau)\Big)\tau^{-m}\beta(\tau)^{-1},\,\,
W_3(\tau)=\|A\|^2,
$$
and

\[
\begin{array}{l}
Z_1(\tau)=2\Big(\sum\limits_{i=1}^2 \eta_i\varphi_i(\tau)\Big)\tau^{-2m}\beta(\tau)^{-2},\vspace{0.1cm}\\ 
Z_2(\tau)=\Big(\sum\limits_{i=1}^2\varphi_i(\tau)^2\Big)\tau^{-2m}\beta(\tau)^{-2}.
\end{array}
\]

If $a_1\eta_1+a_2\eta_2\neq 0$ then
$${\rm{ord}}(W_1(\tau))=-m<1-m\le{\rm{ord}}(2(a_1\varphi_1(\tau)+a_2\varphi_2(\tau)))-m\leq {\rm{ord}}(W_2(\tau)),$$ and ${\rm{ord}}(W_1(\tau))=-m<0={\rm{ord}}(W_3(\tau))$ whence ${\rm{ord}}(W_1(\tau))={\rm{ord}}(W_1(\tau))+W_2(\tau)+W_3(\tau)).$  Observe that $\gamma$ is not a pole of $(\Gamma_A(Y)+\|A\|^2)^{-1}$
otherwise $\gamma$ would be a zero of $\Gamma_A(Y)+\|A\|^2$ and hence $\gamma$ would not be a pole of $\Gamma_A(Y)$.
This implies that ${\rm{ord}}(W_1(\tau))={\rm{ord}}(Z_1(\tau)+Z_2(\tau))$. Hence
$2(a_1\eta_1+a_2\eta_2)\beta(0)^{-1}$ equals the opposite of the coefficient of $\tau^{-m}$ in $Z_1(\tau)+Z_2(\tau)$. Since the latter power series does not depend on $A$, there is a constant $\kappa$ only depending on $\gamma$ such that
\begin{equation}
\label{eqE17apr} a_1\eta_1+a_2\eta_2=\kappa.
\end{equation}
As in the proof of Lemma \ref{lem17apr}, $\gamma$ is also a pole of an other function $\Gamma_T(Y)$, say $\Gamma_D(Y)$. Since (\ref{eqE17apr}) holds true if $A$ is replaced either by $B$ or by $C$, we obtain

\begin{center}
$a_1\eta_1+a_2\eta_2=\kappa,$\\
$d_1\eta_1+d_2\eta_2=\kappa,$\\
$\eta_1^2+\eta_2^2=0.$
\end{center}

Here, $a_1d_2-d_1a_2$ may vanish. In fact, $a_1d_2-d_1a_2=0$ occurs when the points $A,D$ and the origin $O=(0,0)$ are collinear. However, this case can be avoided by a translation which takes $A$ and $B$ to two points which are not collinear with the origin.  Now, from the first two equations,
$$\eta_1=\kappa \frac{d_2-a_2}{a_1d_2-d_1a_2},\quad \eta_2=\kappa \frac{a_1-d_1}{a_1d_2-d_1a_2}.$$
Therefore,
\[
\eta_1^2+\eta_2^2=\kappa^2\frac{{((d_2-a_2)^2+(a_1-d_1)^2)}}{{(a_1d_2-d_1a_2)^2}}.
\]
Since $\eta_1^2+\eta_2^2=0$ this yields $(d_2-a_2)^2+(a_1-d_1)^2=0$ whence $A=D$ follows; but this contradicts the hypothesis that four points $A,B,C,D$ are pairwise distinct.
\end{proof}

\end{document}
